\begin{abstract}
\sloppy
Source attribution of urban PM\(_{2.5}\) from a monitoring network can be posed as a
nonnegative linear inverse problem: proxy emission maps and a wind-driven transport
model give the concentration that each source group produces at each monitor and hour,
and the readings are fitted as a nonnegative combination of these signals plus
background terms.  A fit with a small residual need not determine the contributions: sources aligned with the wind can leave nearly proportional signals at
the sensors, and terms for regional background can absorb a source's signal.  We
present Identifiability-Aware Source Attribution (IASA), which determines from the
transport operator, before fitting, which groups of sources the readings can separate,
and reports each group's contribution with an error bound.  A grouping is identifiable
exactly when the column spaces of its blocks form a direct sum, and the finest such
grouping is unique and computable in polynomial time.  Moreover, the error of a group's
fitted signal is at most the noise in the operator's column space divided by the
group's separation from the other sources.  In controlled experiments on the New Delhi domain,
IASA flags an absorbed source and a nearly proportional pair that the baselines do not
flag, with group-signal errors of \(0.4\)--\(6.0\%\) (\(0.4\)--\(6.2\%\) for nonnegative
least squares after background removal, \(36\)--\(784\%\) for a mass balance without it).
On four observed weeks, the four source groups are separable in every week, and at the
estimated instrument noise level the error bound of the largest group's signal is below
its fitted norm in two of the four weeks.
\end{abstract}
